% ═══════════════════════════════════════════════════════════════════════════════
\section{Theoretische Grundlagen}
\label{sec:grundlagen}
% ═══════════════════════════════════════════════════════════════════════════════

\subsection{Relationale Algebra und SQL}
\label{subsec:rel_algebra}

Die theoretische Basis für die Query-Übersetzung bildet die \emph{relationale
Algebra} nach Codd \cite{codd1970relational}. Alle im UQTL-Standard behandelten
Query-Operatoren lassen sich auf eine endliche Menge algebraischer
Grundoperationen zurückführen.

\begin{definition}[Relationale Grundoperationen]
\label{def:rel_ops}
Gegeben seien Relationen $R$, $S$ über einem Schema $\mathcal{R}$.
Die relationalen Grundoperationen sind:
\begin{align*}
  \sigma_p(R) &\quad \text{Selektion (Filterung) mit Prädikat } p \\
  \pi_A(R)    &\quad \text{Projektion auf Attributmenge } A \subseteq \mathcal{R} \\
  R \times S  &\quad \text{Kartesisches Produkt} \\
  R \cup S    &\quad \text{Vereinigung} \\
  R - S       &\quad \text{Differenz} \\
  R \bowtie_p S &\quad \text{Join mit Prädikat } p \\
  \gamma_{A,F}(R) &\quad \text{Gruppierung nach } A \text{ mit Aggregatfunktionen } F
\end{align*}
\end{definition}

SQL ist eine deklarative Anfragesprache, die über diese algebraischen
Grundoperationen hinausgeht und um Fensterfunktionen, rekursive Abfragen
(CTE), Mengenoperationen sowie datenbankspezifische Erweiterungen ergänzt wird.
ANSI SQL:2016 \cite{ansi_sql_2016} bildet die Ziel-Spezifikation des UQTL-Standards.

\subsection{Query-Abstraktion in modernen Programmiersprachen}
\label{subsec:query_abstraktion}

Alle drei im UQTL-Standard behandelten Sprachen implementieren das Konzept der
\emph{Query Comprehension} oder \emph{monadischen Query-Komposition}:

\subsubsection{Python -- SQLAlchemy Core \& ORM}

Python verfügt über keine native Query-Syntax in der Sprache selbst.
Der De-facto-Standard für relationale Abfragen ist \textbf{SQLAlchemy},
dessen \emph{Core}-Schicht einen Expression Language API bereitstellt.
Abfragen werden als verkettete Methodenaufrufe oder als ORM-Ausdrücke
formuliert \cite{sqlalchemy_docs}:

\begin{lstlisting}[style=pythonstyle, caption={Python SQLAlchemy -- Query-Stilarten},
  label=lst:py_styles]
# Stil 1: SQLAlchemy ORM (Session.query -- Legacy)
users = session.query(User)\
    .filter(User.active == True)\
    .order_by(User.name)\
    .all()

# Stil 2: SQLAlchemy 2.0 Select-Style (modern)
stmt = select(User)\
    .where(User.active == True)\
    .order_by(User.name)
users = session.execute(stmt).scalars().all()

# Stil 3: Pandas / pandasql (fuer Datenpipelines)
import pandasql as ps
users = ps.sqldf("SELECT * FROM users WHERE active=1 ORDER BY name")
\end{lstlisting}

\subsubsection{Java -- Stream API und JPA Criteria API}

Java bietet seit Version 8 die \textbf{Stream API} für funktionale
Datenverarbeitung auf Collections \cite{java_stream_docs}. Für
Datenbankabfragen wird die \textbf{JPA Criteria API} oder Frameworks
wie \textbf{QueryDSL} verwendet:

\begin{lstlisting}[style=javastyle, caption={Java -- Query-Stilarten},
  label=lst:java_styles]
// Stil 1: Stream API (In-Memory-Verarbeitung)
List<User> users = allUsers.stream()
    .filter(u -> u.isActive())
    .sorted(Comparator.comparing(User::getName))
    .collect(Collectors.toList());

// Stil 2: JPA Criteria API (Datenbankabfrage)
CriteriaBuilder cb = em.getCriteriaBuilder();
CriteriaQuery<User> cq = cb.createQuery(User.class);
Root<User> root = cq.from(User.class);
cq.select(root)
  .where(cb.isTrue(root.get("active")))
  .orderBy(cb.asc(root.get("name")));

// Stil 3: QueryDSL (typsicher, fliessend)
List<User> users = new JPAQuery<>(em)
    .select(QUser.user)
    .from(QUser.user)
    .where(QUser.user.active.isTrue())
    .orderBy(QUser.user.name.asc())
    .fetch();
\end{lstlisting}

\subsubsection{C\# -- LINQ (Language Integrated Query)}

C\# bietet mit \textbf{LINQ} eine nativ in die Sprache integrierte
Query-Syntax, die vom Compiler direkt in Methodenaufrufe übersetzt wird.
Entity Framework Core implementiert den \texttt{IQueryable<T>}-Provider
für SQL-Generierung \cite{ef_core_docs,linq_docs}:

\begin{lstlisting}[style=csharpstyle, caption={C\# LINQ -- Query-Stilarten},
  label=lst:cs_styles]
// Stil 1: LINQ Query Syntax (deklarativ, SQL-aehnlich)
var users = from u in context.Users
            where u.Active == true
            orderby u.Name
            select u;

// Stil 2: LINQ Method Syntax (fliessend, Lambda)
var users = context.Users
    .Where(u => u.Active == true)
    .OrderBy(u => u.Name)
    .ToList();

// Stil 3: LINQ mit anonymen Typen (Projektion)
var result = context.Users
    .Where(u => u.Active)
    .Select(u => new { u.Id, u.Name, u.Email });
\end{lstlisting}

\subsection{Bestehende Ansätze und deren Limitierungen}
\label{subsec:bestehend}

Tabelle~\ref{tab:bestehend} fasst bestehende Ansätze zur Query-Abstraktion
und deren Grenzen zusammen.

\begin{table}[H]
\centering
\caption{Bestehende Query-Abstraktionsansätze und ihre Limitierungen}
\label{tab:bestehend}
\renewcommand{\arraystretch}{1.3}
\begin{tabularx}{\textwidth}{lXX}
\toprule
\textbf{Ansatz} & \textbf{Merkmale} & \textbf{Limitierungen} \\
\midrule
Hibernate (Java)
  & Mächtiges ORM, HQL, Criteria API
  & Nur Java; erzeugt oft ineffiziente SQL; kein sprachübergreifender Standard \\
\hline
Entity Framework (C\#)
  & LINQ-Integration, Migrations
  & Nur .NET; SQL-Ausgabe variiert je nach EF-Version \\
\hline
SQLAlchemy (Python)
  & Flexibel, Core + ORM
  & Nur Python; verschiedene API-Generationen; inkonsistente SQL-Ausgabe \\
\hline
JOOQ (Java)
  & Typsicherer SQL-Builder
  & Nur Java; erfordert Code-Generierung; kommerziell für Enterprise \\
\hline
Exposed (Kotlin)
  & Kotlin DSL für SQL
  & Nur JVM/Kotlin; junges Ökosystem \\
\hline
Prisma (Node.js)
  & Typsicherer ORM für TypeScript
  & Nur Node.js/TypeScript; anderes Ökosystem \\
\hline
GraphQL-to-SQL
  & Graphenabfragen über SQL
  & Anderes Paradigma; nicht auf Standard-CRUD ausgerichtet \\
\bottomrule
\end{tabularx}
\end{table}

\begin{warningbox}
\textbf{Kritische Feststellung:}
Kein bestehender Ansatz garantiert sprachübergreifende SQL-Äquivalenz.
Die Lücke liegt nicht in der fehlenden Implementierung, sondern im Fehlen
eines formalen Standards, der Äquivalenzanforderungen und
Normalisierungsregeln verbindlich definiert.
\end{warningbox}

\subsection{Formale Grundbegriffe}
\label{subsec:formale_grundbegriffe}

\begin{definition}[Query-Ausdruck]
\label{def:query_expr}
Ein \emph{Query-Ausdruck} $Q$ ist ein Term einer Operatoralgebra $\mathcal{A}$
über einer Datenbankschema-Instanz $\sigma$. Die \emph{Semantik} von $Q$ ist
die Relation $\llbracket Q \rrbracket_\sigma$, die durch Auswertung von $Q$
in $\sigma$ entsteht.
\end{definition}

\begin{definition}[Semantische Äquivalenz]
\label{def:sem_aequiv}
Zwei Query-Ausdrücke $Q_1$ und $Q_2$ heißen \emph{semantisch äquivalent},
geschrieben $Q_1 \equiv Q_2$, wenn für alle Datenbankinstanzen $\sigma$ gilt:
\[
  \llbracket Q_1 \rrbracket_\sigma = \llbracket Q_2 \rrbracket_\sigma
\]
\end{definition}

\begin{definition}[UQTL-Konformität]
\label{def:konformitaet}
Ein Übersetzungssystem $T$ heißt \emph{UQTL-konform}, wenn für alle
äquivalenten Quellanfragen $Q_{py} \equiv Q_{java} \equiv Q_{cs}$ gilt:
\[
  T(Q_{py}) = T(Q_{java}) = T(Q_{cs}) = S^*
\]
wobei $S^*$ das normalisierte SQL-Statement im UQTL-Normalformat ist.
\end{definition}

Diese Definitionen bilden das formale Fundament des gesamten Standards.

